% Computed scalar profile used by the lower action construction.
% Native TikZ; the preamble loads tikz (no PGFPlots or external data).
% Curves sample Phi(r)=sqrt((1-r)/(8r))+r^2-101/100 directly.
% Main axes: 0.05 <= r <= 1.02, -0.45 <= Phi <= 0.15;
% the graph itself is sampled only for 0.05 <= r < 1.
% Zoom window: 0.965 <= r < 1, -0.012 <= Phi <= 0.007.
% Sampling uses a unit-interval parameter; straight segments avoid spline
% overshoot. PGF floating-point evaluation resolves the shallow right peak.
% Root markers are numerical locations (bisection residual < 5e-15):
% 0.1115973705382644, 0.9786836116269868, 0.9987531118881820.
\begin{figure}[htb]
\centering
\usepgflibrary{fpu}
\pgfmathdeclarefunction{loweractionphi}{1}{%
  \begingroup
  \pgfkeys{/pgf/fpu=true,/pgf/fpu/output format=fixed}%
  \pgfmathparse{sqrt((1-#1)/(8*#1))+#1*#1-1.01}%
  \pgfmathsmuggle\pgfmathresult
  \endgroup
}
\pgfmathdeclarefunction{loweractionphidelta}{1}{%
  \begingroup
  \pgfkeys{/pgf/fpu=true,/pgf/fpu/output format=fixed}%
  \pgfmathparse{sqrt(#1/(8*(1-#1)))-2*#1+#1*#1-.01}%
  \pgfmathsmuggle\pgfmathresult
  \endgroup
}
\begin{tikzpicture}[x=1cm,y=1cm,font=\small]
  % Panel (a): the full profile, with the divergent left tail clipped.
  \begin{scope}
    \draw[line width=.3pt] (0,0) rectangle (7.5,3.5);
    \draw[densely dashed,line width=.35pt] (0,2.625)--(7.5,2.625);
    \foreach \r/\lab in {.1/{0.1},.3/{0.3},.5/{0.5},.7/{0.7},.9/{0.9},1/{1}} {
      \pgfmathsetmacro{\xx}{(\r-.05)/.97*7.5}
      \draw[line width=.3pt] (\xx,0)--(\xx,-.065);
      \node[below,inner sep=2pt] at (\xx,-.065) {$\lab$};
    }
    \foreach \v/\lab in {-.4/{-0.4},-.2/{-0.2},0/{0},.1/{0.1}} {
      \pgfmathsetmacro{\yy}{(\v+.45)/.60*3.5}
      \draw[line width=.3pt] (0,\yy)--(-.065,\yy);
      \node[left,inner sep=2pt] at (-.065,\yy) {$\lab$};
    }
    \begin{scope}
      \clip (0,0) rectangle (7.5,3.5);
      \draw[line width=.75pt]
        plot[domain=0:1,samples=260,variable=\u]
        ({.90*\u/.97*7.5},
         {(loweractionphi(.05+.90*\u)+.45)/.60*3.5})
        plot[domain=0:1,samples=320,variable=\u]
        ({(.90+.04999*\u)/.97*7.5},
         {(loweractionphidelta(.05-.04999*\u)+.45)/.60*3.5});
      \foreach \r in {.1115973705382644,.9786836116269868,.9987531118881820} {
        \fill ({(\r-.05)/.97*7.5},2.625) circle[radius=1.25pt];
      }
    \end{scope}
    \node[above,inner sep=2pt] at (0,3.5) {$\Phi(r)$};
    \node[below,inner sep=2pt] at (7.5,-.42) {$r$};
    \node at (3.75,-.74) {(a) The action profile};
  \end{scope}
  % Panel (b): a separate linear scale resolves both right-hand zeros.
  \begin{scope}[shift={(9.3,0)}]
    \draw[line width=.3pt] (0,0) rectangle (5.2,3.5);
    \draw[densely dashed,line width=.35pt]
      (0,{.012/.019*3.5})--(5.2,{.012/.019*3.5});
    \foreach \r/\lab in {.97/{0.97},.98/{0.98},.99/{0.99},1/{1}} {
      \pgfmathsetmacro{\xx}{(\r-.965)/.035*5.2}
      \draw[line width=.3pt] (\xx,0)--(\xx,-.065);
      \node[below,inner sep=2pt] at (\xx,-.065) {$\lab$};
    }
    \foreach \v/\lab in {-.01/{-0.010},-.005/{-0.005},0/{0},.005/{0.005}} {
      \pgfmathsetmacro{\yy}{(\v+.012)/.019*3.5}
      \draw[line width=.3pt] (0,\yy)--(-.065,\yy);
      \node[left,inner sep=2pt] at (-.065,\yy) {$\lab$};
    }
    \begin{scope}
      \clip (0,0) rectangle (5.2,3.5);
      \draw[line width=.75pt]
        plot[domain=0:1,samples=360,variable=\u]
        ({.03499*\u/.035*5.2},
         {(loweractionphidelta(.035-.03499*\u)+.012)/.019*3.5});
      \foreach \r in {.9786836116269868,.9987531118881820} {
        \fill ({(\r-.965)/.035*5.2},{.012/.019*3.5})
          circle[radius=1.25pt];
      }
    \end{scope}
    \node[above,inner sep=2pt] at (0,3.5) {$\Phi(r)$};
    \node[below,inner sep=2pt] at (5.2,-.42) {$r$};
    \node at (2.6,-.74) {(b) Enlargement near $r=1$};
  \end{scope}
\end{tikzpicture}
\caption{Computed graph of
$\Phi(r)=\sqrt{(1-r)/(8r)}+r^2-101/100$ on $0<r<1$.
Panel~(a) clips the divergent left-end portion as $r\downarrow0$.
The dots mark numerically computed zeros; panel~(b) resolves the two near $r=1$.
The three simple zeros are established analytically in
Lemma~\ref{p4:model}.}
\label{fig:lower-action-profile}
\end{figure}
